\chapter{實驗結果}\label{result}

\section{水之呼吸}
\begin{itemize}
\item 壹之型 水面斬
\item 貳之型 水車
\item 叄之型 流流舞
\item 肆之型 打擊之潮
\item 伍之型 旱天的慈雨
\item 陸之型 扭轉旋渦
\item 柒之型 雫波紋擊刺·曲
\item 玖之型 水流飛沫·亂
\item 拾之型 生生流轉
\item 組合技——陸之型+叄之型——扭轉旋渦·流流
\item 創新技——貳之型·改——橫水車
\item 拾壹之型 凪
\end{itemize}

繪圖範例1：
\begin{figure}[H]

   \centering
        \includegraphics[width=5cm]{../Figures/Result/我就爛}
   \caption{我就爛}
   \label{fig:我就爛}

\end{figure}

% 繪圖範例2：
% \begin{center}
%     \begin{tikzpicture}
%         \begin{axis}[
%             title=Example using the mesh parameter, % 圖標題
%             hide axis,                              % 隱藏座標
%             colormap/cool,                          % 顏色風格
%         ]
%         \addplot3[
%             mesh,                                   % 繪製的三維圖像是網格
%             samples=50,                             % 定義域分割數量
%             domain=-8:8,                            % 定義域
%         ]
%         {sin(deg(sqrt(x^2+y^2)))/sqrt(x^2+y^2)};    % 二元顯式函數
%         \addlegendentry{$\frac{sin(r)}{r}$}         % 添加圖例
%         \end{axis}
%         \end{tikzpicture}
% \end{center}
